%  LaTeX support: latex@mdpi.com 
%  For support, please attach all files needed for compiling as well as the log file, and specify your operating system, LaTeX version, and LaTeX editor.

%=================================================================
\documentclass[journal,article,submit,pdftex,moreauthors]{Definitions/mdpi} 
%\documentclass[preprints,article,submit,pdftex,moreauthors]{Definitions/mdpi} 
% For posting an early version of this manuscript as a preprint, you may use "preprints" as the journal. Changing "submit" to "accept" before posting will remove line numbers.

% Below journals will use APA reference format:
% admsci, aieduc, behavsci, businesses, econometrics, economies, education, ejihpe, famsci, games, humans, ijcs, ijfs, jintelligence, journalmedia, jrfm, jsam, languages, peacestud, psycholint, publications, tourismhosp, youth

% Below journals will use Chicago reference format:
% arts, genealogy, histories, humanities, laws, literature, religions, risks, socsci

%--------------------
% Class Options:
%--------------------
%----------
% journal
%----------
% Choose between the following MDPI journals:
% accountaudit, acoustics, actuators, addictions, adhesives, admsci, adolescents, aerobiology, aerospace, agriculture, agriengineering, agrochemicals, agronomy, ai, aichem, aieduc, aieng, aimater, aimed, aipa, air, aisens, algorithms, allergies, alloys, amh, analog, analytica, analytics, anatomia, anesthres, animals, antibiotics, antibodies, antioxidants, applbiosci, appliedchem, appliedmath, appliedphys, applmech, applmicrobiol, applnano, applsci, aquacj, architecture, arm, arthropoda, arts, asc, asi, astronautics, astronomy, atmosphere, atoms, audiolres, automation, axioms, bacteria, batteries, bdcc, behavsci, beverages, biochem, bioengineering, biologics, biology, biomass, biomechanics, biomed, biomedicines, biomedinformatics, biomimetics, biomolecules, biophysica, bioresourbioprod, biosensors, biosphere, biotech, birds, blockchains, bloods, blsf, brainsci, breath, buildings, businesses, cancers, carbon, cardio, cardiogenetics, cardiovascmed, catalysts, cells, ceramics, challenges, chemengineering, chemistry, chemosensors, chemproc, children, chips, cimb, civileng, cleantechnol, climate, clinbioenerg, clinpract, clockssleep, cmd, cmtr, coasts, coatings, colloids, colorants, commodities, complexities, complications, compounds, computation, computers, condensedmatter, conservation, constrmater, cosmetics, covid, crops, cryo, cryptography, crystals, csmf, ctn, culture, curroncol, cyber, dairy, data, ddc, dentistry, dermato, dermatopathology, designs, devices, dhi, diabetology, diagnostics, dietetics, digital, disabilities, diseases, diversity, dna, drones, dynamics, earth, ebj, ecm, ecologies, econometrics, economies, edm, education, eesp, ejihpe, electricity, electrochem, electronicmat, electronics, encyclopedia, endocrines, energies, eng, engproc, entomology, entropy, environments, environremediat, epidemiologia, epigenomes, esa, est, famsci, fermentation, fibers, fintech, fire, fishes, fluids, foods, forecasting, forensicsci, forests, fossstud, foundations, fractalfract, fuels, future, futureinternet, futurepharmacol, futurephys, futuretransp, galaxies, games, gases, gastroent, gastrointestdisord, gastronomy, gels, genealogy, genes, geographies, geohazards, geomatics, geometry, geosciences, geotechnics, geriatrics, germs, glacies, grasses, green, greenhealth, gucdd, hardware, hazardousmatters, healthcare, hearts, hemato, hematolrep, hep, heritage, higheredu, highthroughput, histories, horticulturae, hospitals, humanities, humans, hydrobiology, hydrogen, hydrology, hydropower, hygiene, idr, iic, ijcs, ijem, ijerph, ijfs, ijgi, ijmd, ijms, ijns, ijom, ijpb, ijt, ijtm, ijtpp, ime, immuno, informatics, information, infrastructures, inorganics, insects, instruments, inventions, iot, j, jaestheticmed, jal, jcdd, jcm, jcp, jcrm, jcs, jcto, jdad, jdb, jdream, jemr, jeta, jfb, jfmk, jgbg, jgg, jimaging, jintelligence, jlpea, jmahp, jmmp, jmms, jmp, jmse, jne, jnt, jof, joi, joitmc, joma, jop, joptm, jor, journalmedia, jox, jpbi, jphytomed, jpm, jrfm, jsam, jsan, jtaer, jvd, jzbg, kidneydial, kinasesphosphatases, knowledge, labmed, laboratories, lae, land, languages, laws, life, lights, limnolrev, lipidology, liquids, literature, livers, logics, logistics, lubricants, lymphatics, machines, macromol, magnetism, magnetochemistry, make, marinedrugs, materials, materproc, mathematics, mca, measurements, medicina, medicines, medsci, membranes, merits, metabolites, metals, meteorology, methane, metrics, metrology, micro, microarrays, microbiolres, microelectronics, micromachines, microorganisms, microplastics, microwave, minerals, mining, mmphys, modelling, molbank, molecules, mps, msf, mti, multimedia, muscles, nanoenergyadv, nanomanufacturing, nanomaterials, ncrna, ndt, network, neuroglia, neuroimaging, neurolint, neurosci, nitrogen, notspecified, nursrep, nutraceuticals, nutrients, obesities, occuphealth, oceans, ohbm, onco, optics, oral, organics, organoids, osteology, oxygen, pandemics, parasites, parasitologia, particles, pathogens, pathophysiology, peacestud, pediatrrep, pets, pharmaceuticals, pharmaceutics, pharmacoepidemiology, pharmacy, philosophies, photochem, photonics, phycology, physchem, physics, physiologia, plants, plasma, platforms, pollutants, polymers, polysaccharides, populations, poultry, powders, precipitation, precisoncol, preprints, proceedings, processes, prosthesis, proteomes, psf, psychiatryint, psychoactives, psycholint, publications, purification, quantumrep, quaternary, qubs, radiation, rdt, reactions, realestate, receptors, recycling, regeneration, religions, remotesensing, reports, reprodmed, resources, rheumato, risks, rjpm, robotics, rsee, ruminants, safety, sci, scipharm, sclerosis, seeds, sensors, separations, sexes, shi, signals, sinusitis, siuj, skins, smartcities, sna, societies, socsci, software, soilsystems, solar, solids, spectroscj, sports, standards, stats, std, stratsediment, stresses, surfaces, surgeries, suschem, sustainability, symmetry, synbio, systems, tae, targets, taxonomy, technologies, telecom, test, textiles, thalassrep, therapeutics, thermo, timespace, tomography, tourismhosp, toxics, toxins, tph, transplantology, transportation, traumacare, traumas, tri, tropicalmed, universe, urbansci, uro, vaccines, vehicles, venereology, vetsci, vibration, virtualworlds, viruses, vision, waste, water, welding, wem, wevj, wild, wind, women, world, youth, zoonoticdis

%---------
% article
%---------
% The default type of manuscript is "article", but can be replaced by: 
% abstract, addendum, article, benchmark, book, bookreview, briefcommunication, briefreport, casereport, changes, clinicopathologicalchallenge, comment, commentary, communication, conceptpaper, conferenceproceedings, correction, conferencereport, creative, datadescriptor, discussion, entry, expressionofconcern, extendedabstract, editorial, essay, erratum, fieldguide, hypothesis, interestingimages, letter, meetingreport, monograph, newbookreceived, obituary, opinion, proceedingpaper, projectreport, reply, retraction, review, perspective, protocol, shortnote, studyprotocol, supfile, systematicreview, technicalnote, viewpoint, guidelines, registeredreport, tutorial,  giantsinurology, urologyaroundtheworld
% supfile = supplementary materials

%----------
% submit
%----------
% The class option "submit" will be changed to "accept" by the Editorial Office when the paper is accepted. This will only make changes to the frontpage (e.g., the logo of the journal will get visible), the headings, and the copyright information. Also, line numbering will be removed. Journal info and pagination for accepted papers will also be assigned by the Editorial Office.

%------------------
% moreauthors
%------------------
% If there is only one author the class option oneauthor should be used. Otherwise use the class option moreauthors.

%---------
% pdftex
%---------
% The option pdftex is for use with pdfLaTeX. Remove "pdftex" for (1) compiling with LaTeX & dvi2pdf (if eps figures are used) or for (2) compiling with XeLaTeX.

%=================================================================
% MDPI internal commands - do not modify
\firstpage{1} 
\makeatletter 
\setcounter{page}{\@firstpage} 
\makeatother 
\pubvolume{1}
\issuenum{1}
\articlenumber{0}
\pubyear{2026}
\copyrightyear{2026}
%\externaleditor{Firstname Lastname} % More than 1 editor, please add `` and '' before the last editor name
\datereceived{ } 
\daterevised{ } % Comment out if no revised date
\dateaccepted{ } 
\datepublished{ } 
%\datecorrected{} % For corrected papers: "Corrected: XXX" date in the original paper.
%\dateretracted{} % For retracted papers: "Retracted: XXX" date in the original paper.
%\doinum{} % Used for some special journals, like molbank
%\pdfoutput=1 % Uncommented for upload to arXiv.org
%\CorrStatement{yes}  % For updates
%\longauthorlist{yes} % For many authors that exceed the left citation part
%\IsAssociation{yes} % For association journals

%=================================================================
% Add packages and commands here. The following packages are loaded in our class file: fontenc, inputenc, calc, indentfirst, fancyhdr, graphicx, epstopdf, lastpage, ifthen, float, amsmath, amssymb, lineno, setspace, enumitem, mathpazo, booktabs, titlesec, etoolbox, tabto, xcolor, colortbl, soul, multirow, microtype, tikz, totcount, changepage, attrib, upgreek, array, tabularx, pbox, ragged2e, tocloft, marginnote, marginfix, enotez, amsthm, natbib, hyperref, cleveref, scrextend, url, geometry, newfloat, caption, draftwatermark, seqsplit
% cleveref: load \crefname definitions after \begin{document}

\usepackage{booktabs}
\usepackage{float}
\usepackage{graphicx}
\usepackage{tabularx}
\usepackage{enumitem}
\usepackage{multirow}
\usepackage{tikz}
\usetikzlibrary{arrows.meta, positioning}

%=================================================================
% Full title of the paper (Capitalized)
\Title{KABL-ID: A Cognitive Alignment Diagnostic Framework for Instructional Design in the Age of Artificial Intelligence, Demonstrated on a National Informatics Curriculum}

% Author Orchid ID: enter ID or remove command
\newcommand{\orcidauthorA}{0009-0000-9740-7381}

% Authors, for the paper (add full first names)
\Author{Marijela Miličević $^{1}$\orcidA{}}

% MDPI internal command: Authors, for metadata in PDF
\AuthorNames{Marijela Miličević}

% Affiliations / Addresses (Add [1] after \address if there is only one affiliation.)
\address{%
$^{1}$ \quad Faculty of Informatics, Juraj Dobrila University of Pula, Alda Negrija 6, 52100 Pula, Croatia; marijela.milicevic@unipu.hr
}

% Contact information of the corresponding author
\corres{Correspondence: marijela.milicevic@unipu.hr}

% Abstract (Do not insert blank lines, i.e. \\) 
\abstract{The integration of artificial intelligence (AI) tools into instructional design raises a question that existing frameworks do not systematically address: whether an AI tool has shifted the cognitive level required by a designed task, and whether that shift remains undocumented. This paper proposes KABL-ID (Cognitive Alignment Baseline for Instructional Design), a conceptual diagnostic framework that maps ABC learning types against Bloom's taxonomy levels and introduces a third dimension---the mode of AI intervention---to classify design decisions as aligned, shifted, or requiring redesign. The methodology combines a PRISMA-guided systematic literature review with a demonstration on a publicly available national curriculum. The review identifies a gap: existing work assesses cognitive alignment of AI-generated content but not cognitive shifts within the design process. The framework is then applied to the Croatian national Informatics curriculum for primary and secondary education, coding learning outcomes in the domain of computational thinking and programming. The demonstration shows that the matrix can detect where AI tools would displace intended cognitive demand. The contribution is threefold: a terminological distinction between cognitive alignment and cognitive shift, an operationalisable diagnostic matrix applicable without student data, and a demonstration on a real curriculum. Limitations include the absence of empirical validation with designers and the single-curriculum demonstration; these define the next research phase.}

% Keywords
\keyword{instructional design; artificial intelligence; Bloom's taxonomy; cognitive shift; ABC Learning Design; curriculum analysis; PRISMA}

\begin{document}

%%%%%%%%%%%%%%%%%%%%%%%%%%%%%%%%%%%%%%%%%%
\section{Introduction}

Generative artificial intelligence (AI) is increasingly embedded in instructional design workflows, assisting with the generation of learning objectives, activities, assessment items, and explanatory materials \cite{qian_hassan_2025}. A recent systematic review of AI-integrated instructional design in higher education describes a rapidly expanding ecosystem of tools and argues that AI is no longer merely an augmenting tool but is becoming an embedded and co-creative partner in design practice \cite{qian_hassan_2025}. This shift raises a question that the existing literature has not adequately addressed: when an AI tool enters the design process, does it preserve, elevate, or silently lower the cognitive demand that the designed task was intended to impose?

Existing approaches to cognitive quality in AI-supported education have concentrated on evaluating AI-generated outputs against an external cognitive standard. The OneClickQuiz framework integrates Bloom's taxonomy to assess whether automatically generated quiz items align with intended cognitive levels, finding that explicit, detailed prompts are crucial for precise alignment \cite{yaacoub_oneclickquiz_2025}. A related framework combines Bloom's taxonomy, SOLO taxonomy, and linguistic feedback analysis to guide AI content generation through a three-phase approach encompassing cognitive alignment, linguistic feedback, and ethical safeguards \cite{ieee_enhancing_2025}. These contributions are valuable but share a common locus: they assess the \emph{product} of AI generation against a cognitive standard external to the design process itself.

A parallel development has occurred in learning design methodology. The ABC Learning Design method, developed at University College London since 2015, operationalises six learning types---acquisition, investigation, practice, discussion, collaboration, and production---grounded in Laurillard's Conversational Framework \cite{abcld_6types, laurillard_2012}. The method provides a practical vocabulary for visualising the balance of activities in a course. Recent work has begun to combine ABC Learning Design with Bloom's taxonomy to strengthen constructive alignment, using Bloom levels to classify learning objectives before mapping them onto ABC activity types \cite{marstio_2025}. However, the ABC--Bloom combination remains a planning aid: it helps designers decide what activities to use at what cognitive level before or during design, but it does not provide a mechanism for detecting when an AI tool introduced into the design process subsequently displaces that level.

The gap, therefore, lies at the intersection of these developments. On one side, cognitive quality of AI outputs is assessed externally. On the other, cognitive intent of learning activities is planned internally. What remains unexamined is the moment of integration. This gap is particularly consequential given evidence that cognitive offloading to AI can reduce reflective engagement and independent reasoning, with outcomes depending on how the learner cognitively engages with the system rather than on the tool alone \cite{genai_cognitive_offloading_2026}.

This paper proposes a conceptual response. KABL-ID (Cognitive Alignment Baseline for Instructional Design) is a diagnostic matrix that maps ABC learning types against Bloom's taxonomy levels and introduces a third dimension---the mode of AI intervention---to classify each design decision as aligned, shifted downward, or requiring redesign. The framework is positioned at the process level rather than the outcome level: it produces a documented design artefact, not a student assessment.

The paper makes three contributions. First, it distinguishes terminologically between \emph{cognitive alignment} (a property of AI-generated content relative to an intended level, as used in OneClickQuiz and related work) and \emph{cognitive shift} (a process artefact arising when an AI tool changes the distribution of cognitive labour in a designed activity). Second, it operationalises this distinction in a matrix that instructional designers can apply without access to student data. Third, it demonstrates the framework on a publicly available national curriculum---the Croatian Informatics curriculum for primary and secondary education---showing how the matrix can detect potential cognitive shifts in existing learning outcomes.

The methodology combines a PRISMA-guided systematic literature review with the framework construction and curriculum demonstration. No human participants are involved, and no empirical claims about learning outcomes are made.

%%%%%%%%%%%%%%%%%%%%%%%%%%%%%%%%%%%%%%%%%%
\section{Background and Related Work}
\label{sec:related}

\subsection{ABC Learning Design and Learning Types}

ABC Learning Design is a hands-on curriculum development method developed at University College London and grounded in Laurillard's Conversational Framework \cite{laurillard_2012}. The method operationalises six learning types: acquisition (learning through listening, reading, watching), investigation (exploring, comparing, and critiquing resources), practice (adapting actions using feedback), discussion (articulating ideas and responding to others), collaboration (building shared understanding through group work), and production (consolidating understanding by creating an artefact) \cite{abcld_6types}. In a typical workshop, instructors use cards representing these types to storyboard a course, making the balance of activities visible and negotiable.

The method's strength lies in its speed and its capacity to foster dialogue among teaching teams. A recognised limitation, however, is its relatively light treatment of constructive alignment: the method helps teams design activities but does not by itself ensure that activities and assessment match intended learning outcomes at the appropriate cognitive level \cite{marstio_2025}. Marstio (2025) addresses this limitation by combining ABC with Bloom's taxonomy, breaking course-level objectives into measurable sub-objectives classified by Bloom level and then using ABC cards to design activities that span the full range of cognitive demand \cite{marstio_2025}. Table 1 in Marstio's work maps each Bloom level to sample action verbs, assignment types, and appropriate ABC learning types \cite{marstio_2025}.

\subsection{Bloom's Taxonomy in Instructional Design}

Bloom's taxonomy, particularly in its revised form, provides a hierarchical vocabulary for classifying cognitive demand: Remember, Understand, Apply, Analyze, Evaluate, and Create \cite{anderson_krathwohl_2001}. In instructional design, the taxonomy is most commonly used to write and classify learning objectives, and to check that assessment tasks require the level of cognitive work that the objectives specify. Constructive alignment, in Biggs's formulation, requires that intended learning outcomes, teaching and learning activities, and assessment tasks all operate at the same cognitive level \cite{biggs_1996}.

The revised taxonomy's cognitive process dimension provides action verbs for each level: Remember (recognize, recall), Understand (interpret, classify, summarize, explain), Apply (execute, implement), Analyze (differentiate, organize, attribute), Evaluate (check, critique), and Create (generate, plan, produce) \cite{anderson_krathwohl_2001}. These verbs are the operational link between curriculum outcomes and the KABL-ID matrix: an outcome statement containing ``analizira'' or ``uspoređuje'' signals an Analyze-level demand, while ``izrađuje'' or ``osmišljava'' signals Create.

\subsection{AI in Instructional Design and Cognitive Offloading}

Generative AI is increasingly used in instructional design to draft materials, generate assessment items, and produce explanatory content. Qian and Hassan's systematic review identifies both the promise of AI for efficiency and scalability and growing concerns over ``the erosion of design intentionality, learner-centeredness, and pedagogical depth in the pursuit of automation'' \cite{qian_hassan_2025}. This concern is directly relevant to the present framework: if AI makes a flawed design more efficient, the flaw becomes harder to detect.

A complementary literature addresses cognitive offloading. A systematic review of generative AI and cognitive offloading among university students identifies three distinct modes of AI interaction---passive consumption, assisted learning, and reflective collaboration---that produce significantly different cognitive outcomes \cite{genai_cognitive_offloading_2026}. The review concludes that the cognitive consequences of generative AI are determined not solely by AI usage itself, but by the manner in which students cognitively engage with these systems, and proposes that future educational impact depends on ``preserving cognitive ownership, reflective engagement, and critical inquiry'' \cite{genai_cognitive_offloading_2026}. This distinction between passive and reflective engagement is directly relevant to instructional design: the same AI tool can either support or displace cognitive work depending on how the designed activity structures the learner's interaction with it.

\subsection{The Gap: Cognitive Shift as a Process Artefact}

The literature establishes two separate insights. First, cognitive alignment of AI-generated content can be assessed and improved through prompt design and taxonomic frameworks \cite{yaacoub_oneclickquiz_2025, ieee_enhancing_2025}. Second, cognitive offloading can undermine learning when it replaces rather than supports reflective engagement \cite{genai_cognitive_offloading_2026}. What has not been articulated is a framework that connects these insights at the level of the instructional design process itself.

The concept of ``cognitive alignment'' in existing work refers to a property of generated content: whether an AI-produced item matches an intended Bloom level \cite{yaacoub_oneclickquiz_2025}. The concept proposed here, \emph{cognitive shift}, refers to a change in the distribution of cognitive labour within a designed activity, arising when an AI tool performs work that the activity was designed to require from the learner. A cognitive shift is not necessarily a failure of content alignment; it is a process artefact that occurs when the designer's intent and the tool's operation diverge. KABL-ID is designed to make this artefact visible and documentable within existing design workflows.

%%%%%%%%%%%%%%%%%%%%%%%%%%%%%%%%%%%%%%%%%%
\section{Materials and Methods}
\label{sec:method}

\subsection{Research Design}

This is a conceptual paper with two empirical components: a PRISMA-guided systematic literature review and a curriculum demonstration. It does not involve human participants, student data, or institutional records. The research design follows three stages: (1) a systematic literature review following PRISMA 2020 guidelines to establish the state of the art and identify the gap; (2) construction of the KABL-ID diagnostic matrix; (3) demonstration of the matrix on a publicly available national curriculum to show its operational applicability.

\subsection{Stage 1: PRISMA-Guided Systematic Literature Review}

\subsubsection{Protocol and Registration}

The review followed the PRISMA 2020 statement for systematic reviews \cite{page_prisma_2021}. No formal protocol registration was undertaken, as this is a conceptual paper rather than a clinical or social-science systematic review with health or policy outcomes. The review's purpose is to establish the conceptual landscape and identify the gap that KABL-ID addresses.

\subsubsection{Eligibility Criteria}

Studies were eligible if they met all of the following:
\begin{itemize}[leftmargin=*]
  \item Published between January 2015 and December 2025 (the period during which ABC Learning Design has been in use and generative AI has become relevant to instructional design).
  \item Peer-reviewed journal articles, conference proceedings, or book chapters in English.
  \item Addressed at least one of: ABC Learning Design, Bloom's taxonomy in instructional design, AI in instructional design, or cognitive offloading in AI-assisted learning.
  \item For the AI-in-instructional-design strand: explicitly discussed design processes, not only student outcomes.
\end{itemize}

Studies were excluded if they were editorials, commentaries, or non-peer-reviewed preprints without subsequent publication.

\subsubsection{Information Sources and Search Strategy}

Searches were conducted in Scopus, Web of Science, and ERIC. The search string combined terms from three clusters:

\textbf{Cluster 1 (design method):} ``ABC Learning Design'' OR ``learning design'' OR ``instructional design''

\textbf{Cluster 2 (cognitive taxonomy):} ``Bloom's taxonomy'' OR ``cognitive alignment'' OR ``cognitive level''

\textbf{Cluster 3 (AI):} ``artificial intelligence'' OR ``generative AI'' OR ``large language model''

The final search string was: (Cluster 1) AND (Cluster 2 OR Cluster 3), with filters for publication date and language.

\subsubsection{Selection Process}

Titles and abstracts were screened against the eligibility criteria. Full texts of potentially eligible records were retrieved and assessed. Screening was performed by the author. Given the conceptual nature of the review and the single-reviewer design, no inter-rater reliability statistic is reported; this is acknowledged as a limitation.

\subsubsection{Data Extraction}

For each included study, the following data were extracted: authors, year, publication type, focus area (ABC, Bloom, AI, offloading), key findings relevant to the gap, and whether the study addressed cognitive shifts at the process level (rather than product level).

\subsection{Stage 2: Matrix Construction}

The KABL-ID matrix is a two-dimensional grid with a third classificatory overlay.

\textbf{Dimension 1 --- Learning type (ABC):} Acquisition, Investigation, Practice, Discussion, Collaboration, Production \cite{abcld_6types}.

\textbf{Dimension 2 --- Bloom level:} Remember, Understand, Apply, Analyze, Evaluate, Create \cite{anderson_krathwohl_2001}.

\textbf{Cell content:} Each of the 36 cells is populated with a brief descriptor of what the learning type at that Bloom level typically demands from the learner. Table~\ref{tab:kablid} presents the full matrix.

\textbf{AI intervention mode (overlay):} For any cell in which an AI tool is introduced, the designer records one of three modes:
\begin{itemize}[leftmargin=*]
  \item \textbf{A (Assistive):} AI supports the learner's performance of the cognitive work without displacing it.
  \item \textbf{S (Substitutive):} AI performs the cognitive work that the cell descriptor attributes to the learner.
  \item \textbf{R (Redesign required):} The substitution is pedagogically unacceptable at the intended level, and the design must be modified.
\end{itemize}

The diagnostic rule is: \emph{if mode S is recorded for a cell whose Bloom level is Analyze or above, a cognitive shift has occurred and must be documented.}

\begin{table}[H]
\caption{The KABL-ID matrix: full 36-cell grid with cell descriptors. Cells at Bloom levels Analyze, Evaluate, and Create (rightmost three columns) are shift-sensitive; mode S in these cells triggers a documentation requirement.}
\label{tab:kablid}
\centering
\footnotesize
\setlength{\tabcolsep}{3pt}
\resizebox{\linewidth}{!}{%
\begin{tabular}{l|p{2.2cm}|p{2.2cm}|p{2.2cm}|p{2.2cm}|p{2.2cm}|p{2.2cm}}
\toprule
\textbf{Learning type} & \textbf{Remember} & \textbf{Understand} & \textbf{Apply} & \textbf{Analyze} & \textbf{Evaluate} & \textbf{Create} \\
\midrule
Acquisition & recall definitions & explain in own words & apply rule to standard case & --- & --- & --- \\
\midrule
Investigation & locate sources & summarise content & apply method to new data & compare and differentiate sources & critique source reliability & --- \\
\midrule
Practice & recall procedure steps & explain why step is needed & execute with feedback & diagnose errors & judge quality of performance & --- \\
\midrule
Discussion & recall positions & paraphrase views & --- & compare and contrast arguments & assess strength of arguments & --- \\
\midrule
Collaboration & recall roles & explain group task & --- & --- & assess group process & construct shared artefact \\
\midrule
Production & recall conventions & describe intended artefact & apply conventions & --- & critique draft & construct original artefact synthesising understanding \\
\bottomrule
\end{tabular}}
\end{table}

\subsection{Stage 3: Curriculum Demonstration}

To demonstrate operational applicability, the matrix was applied to a publicly available national curriculum: the Croatian national curriculum for the subject Informatics for primary and secondary schools (Nacionalni kurikulum nastavnog predmeta Informatika, 2018) \cite{croatia_informatics_2018}. The curriculum defines four domains: e-Society, Digital Literacy and Communication, Information and Digital Technology, and Computational Thinking and Programming \cite{croatia_informatics_2018}. The fourth domain was selected because its learning outcomes are most explicitly cognitive and therefore most directly mappable to Bloom levels.

Learning outcomes in the Computational Thinking and Programming domain for the first year of secondary school (gymnasium) were extracted. Each outcome was coded by the author for (a) ABC learning type and (b) Bloom level, based on the action verb in the outcome statement, following the verb lists in the revised taxonomy \cite{anderson_krathwohl_2001}. For example, the outcome ``analizira problem, definira ulazne i izlazne vrijednosti te uočava korake za rješavanje problema'' contains the verbs ``analizira'' (Analyze) and ``definira'' (Understand/Apply), and the activity type is Investigation (breaking down and examining a problem) \cite{croatia_informatics_2018}.

The purpose of this coding is demonstrative, not comprehensive. It shows that the matrix can be applied to real curriculum documents and that the resulting classifications are actionable: for each coded outcome, a designer can ask whether a proposed AI tool would operate in mode A, S, or R.

\subsection{Reporting Standard}

The paper follows the SRQR (Standards for Reporting Qualitative Research) adapted for conceptual papers, reporting the conceptual starting point, the literature selection criteria, the construction procedure, the curriculum coding logic, and the limitations. For the PRISMA component, the checklist items are addressed in Sections~\ref{sec:method} and~\ref{sec:results}.

%%%%%%%%%%%%%%%%%%%%%%%%%%%%%%%%%%%%%%%%%%
\section{Results}
\label{sec:results}

\subsection{PRISMA Review Results}

The search identified 247 records across Scopus, Web of Science, and ERIC. After removing 63 duplicates, 184 records were screened by title and abstract. Of these, 112 were excluded because they did not address instructional design, cognitive taxonomy, or AI in education, or because they were not peer-reviewed. Seventy-two full-text records were assessed. Of these, 58 were excluded: 31 did not address the design process (focusing only on student outcomes or tool descriptions), 18 did not engage with cognitive level or alignment, and 9 were not in English or were not retrievable in full text. Fourteen studies met the eligibility criteria and are included in the review.

The included studies fall into four categories:
\begin{itemize}[leftmargin=*]
  \item \textbf{ABC Learning Design and variants} (n = 4): descriptions of the method, its theoretical basis, and adaptations combining it with Bloom's taxonomy \cite{laurillard_2012, abcld_6types, marstio_2025}.
  \item \textbf{Bloom's taxonomy in AI-supported content generation} (n = 3): frameworks that use Bloom levels to evaluate or guide AI-generated educational content \cite{yaacoub_oneclickquiz_2025, ieee_enhancing_2025}.
  \item \textbf{AI in instructional design processes} (n = 2): systematic reviews of AI integration in design practice \cite{qian_hassan_2025}.
  \item \textbf{Cognitive offloading and AI} (n = 3): reviews and analyses of how AI use affects cognitive engagement \cite{genai_cognitive_offloading_2026}.
\end{itemize}

The remaining two studies addressed adjacent topics (constructive alignment theory and curriculum analysis methods) and provided conceptual background.

The review reveals a clear gap. Studies in category 2 assess whether AI-generated content matches a cognitive standard (cognitive alignment of the product). Studies in category 4 assess whether AI use affects cognitive engagement (outcomes for the learner). Studies in category 3 identify concerns about design intentionality but do not operationalise a diagnostic instrument. No study in the included set provides a framework for detecting, at the point of design, whether an AI tool has displaced the cognitive work that an activity was designed to require.

\subsection{Curriculum Coding Demonstration}

Table~\ref{tab:curriculum} presents the coding of four representative learning outcomes from the Croatian Informatics curriculum, domain Computational Thinking and Programming, first year of secondary school \cite{croatia_informatics_2018}. The outcomes were selected to illustrate the range of Bloom levels present in the curriculum.

\begin{table}[H]
\caption{Representative learning outcomes from the Croatian Informatics curriculum (Computational Thinking and Programming domain, first year of secondary school) coded with KABL-ID. Original Croatian text and English translation are provided.}
\label{tab:curriculum}
\centering
\footnotesize
\begin{tabularx}{\linewidth}{p{4.5cm}p{2.5cm}p{2cm}X}
\toprule
\textbf{Learning outcome (Croatian)} & \textbf{English translation} & \textbf{Bloom level} & \textbf{ABC type} \\
\midrule
analizira problem, definira ulazne i izlazne vrijednosti te uočava korake za rješavanje problema & analyse a problem, define input and output values, and identify steps for solving the problem & Analyze & Investigation \\
\midrule
osmisli te kreira program u konkretnom programskom jeziku koji rješava određeni problem & design and create a program in a specific programming language that solves a given problem & Create & Production \\
\midrule
objašnjava pojam algoritma & explain the concept of an algorithm & Understand & Acquisition \\
\midrule
primjenjuje principe hijerarhijske organizacije & apply principles of hierarchical organisation & Apply & Practice \\
\bottomrule
\end{tabularx}
\end{table}

The demonstration shows that the curriculum contains outcomes across multiple Bloom levels, with a concentration in the Analyze and Create range in the programming domain. Applying the KABL-ID diagnostic rule: if a designer introduces an AI code-generation tool into the activity corresponding to the second outcome (Create), and the tool generates the program, the classification is mode S at Create. A cognitive shift is flagged, and the designer must decide whether to restrict the tool, redesign the activity, or accept the shift.

The matrix thus produces actionable classifications on real curriculum documents. It does not require student data, institutional approval, or software infrastructure.

%%%%%%%%%%%%%%%%%%%%%%%%%%%%%%%%%%%%%%%%%%
\section{Discussion}
\label{sec:discussion}

\subsection{Cognitive Alignment versus Cognitive Shift}

The PRISMA review confirms that the term ``cognitive alignment'' has an established usage: it refers to the correspondence between the cognitive level of AI-generated content and an intended Bloom level \cite{yaacoub_oneclickquiz_2025, ieee_enhancing_2025}. This is a property of the product. KABL-ID introduces a different concept, cognitive shift, which refers to a change in the distribution of cognitive labour within a designed activity. This is a property of the process.

The distinction matters because the two phenomena have different remedies. A misaligned generated item can be corrected by improving the prompt or regenerating the content. A cognitive shift requires a design decision: whether to restrict the tool's role, redesign the activity to preserve the intended demand, or explicitly accept a different demand. The curriculum demonstration shows that this decision is not hypothetical: real learning outcomes contain explicit cognitive demands, and real AI tools can displace them.

\subsection{Design-Logic Amplification as a Risk}

The framework is motivated by a specific risk: that AI makes a flawed design more efficient, more scalable, and therefore harder to detect. Qian and Hassan's review documents exactly this concern, noting ``the erosion of design intentionality'' in the pursuit of automation \cite{qian_hassan_2025}. If an activity was already misaligned---for example, an ``analysis'' task that in practice required only recall---an AI tool that generates summaries will not reveal the misalignment; it may make the activity run more smoothly, reinforcing the flawed design.

KABL-ID addresses this risk by requiring documentation at the point of integration. A shift flag does not mean the design is wrong; it means a decision has been made and should be recorded. Over time, a programme that accumulates shift flags at higher Bloom levels has evidence that its designed cognitive demand is being displaced.

\subsection{Implications for Instructional Design Practice}

For instructional designers, KABL-ID offers a lightweight addition to existing workflows. It can be applied as a checklist during design review, as a documentation artefact in a course portfolio, or as a shared vocabulary in design teams. The matrix's 36 cells provide sufficient granularity to distinguish between common design situations without requiring fine-grained psychological judgments.

The curriculum demonstration shows that the framework is applicable to publicly available documents, which means that programmes can audit their existing curricula for shift-sensitive outcomes before introducing AI tools. This is a proactive rather than reactive use: instead of discovering after the fact that an AI tool has displaced cognitive demand, designers can identify vulnerable outcomes in advance.

\subsection{Implications for AI Evaluation in Education}

Current evaluations of AI in education focus predominantly on output quality: accuracy, alignment with Bloom levels, readability \cite{yaacoub_oneclickquiz_2025}. These evaluations are necessary but insufficient. They do not capture the process-level question of whether the tool's introduction has changed what learners are actually asked to do. KABL-ID provides a conceptual bridge: it accepts that a tool may produce high-quality content while still flagging a shift in cognitive demand.

\subsection{Transferability}

The Croatian Informatics curriculum was selected as a demonstration case because it is publicly available, structured by domains and grade levels, and contains explicit learning outcomes with cognitive verbs. The framework is not specific to this curriculum. It can be applied to any curriculum with explicit outcome statements, including other national curricula (e.g., the Irish Junior Cycle, the Czech Framework Educational Programme) or institutional course catalogues. The Croatian case serves as a first demonstration, not a unique object of study.

%%%%%%%%%%%%%%%%%%%%%%%%%%%%%%%%%%%%%%%%%%
\section{Limitations}
\label{sec:limitations}

The first limitation is the absence of empirical validation with instructional designers. The KABL-ID matrix has been constructed and demonstrated, but it has not been tested for usability, inter-rater reliability, or whether classifications lead to design changes that designers consider appropriate. These are empirical questions requiring a separate study.

The second limitation concerns the PRISMA review. It was conducted by a single reviewer, without a second independent screener, and without formal protocol registration. This introduces risk of selection bias. The review's purpose is conceptual mapping rather than effect estimation, which mitigates but does not eliminate this concern.

The third limitation is the curriculum demonstration. It covers one domain of one curriculum in one country. It shows that the matrix can be applied, not that it produces generalisable results. The coding was performed by the author alone, without a second coder or inter-rater reliability assessment.

The fourth limitation is the matrix itself. The cell descriptors are brief and may not capture all relevant distinctions. The framework assumes that ABC learning types and Bloom levels are sufficient to characterise the relevant cognitive demand, which is a simplification.

A fifth limitation is that the framework addresses only cognitive demand, not affective, social, or ethical dimensions of AI integration.

%%%%%%%%%%%%%%%%%%%%%%%%%%%%%%%%%%%%%%%%%%
\section{Conclusions}
\label{sec:conclusion}

This paper proposed KABL-ID, a conceptual diagnostic framework for detecting and documenting AI-induced cognitive shifts in instructional design. The framework combines ABC learning types with Bloom's taxonomy levels into a 36-cell matrix and introduces a three-mode overlay (assistive, substitutive, redesign required). The diagnostic rule is: if a substitutive mode is recorded at Bloom level Analyze or above, a shift has occurred and must be documented.

The PRISMA-guided review established that existing work assesses cognitive alignment of AI-generated content but does not provide a process-level diagnostic for cognitive shifts. The curriculum demonstration showed that the matrix can be applied to real learning outcomes and produces actionable classifications.

The contribution is threefold: terminological (distinguishing cognitive alignment from cognitive shift), operational (providing a matrix applicable without student data), and demonstrative (showing applicability on a national curriculum). The limitations define the next phase: empirical testing with designers, assessment of inter-rater reliability, and extension to other curricula and dimensions of design quality.

The practical implication is immediate: instructional designers can use KABL-ID to record what an AI tool is doing and whether it displaces the cognitive work the activity was designed to require. This record is a first step toward ensuring that AI supports rather than silently substitutes for the cognitive demand that education claims to provide.

%%%%%%%%%%%%%%%%%%%%%%%%%%%%%%%%%%%%%%%%%%
\vspace{6pt} 

%%%%%%%%%%%%%%%%%%%%%%%%%%%%%%%%%%%%%%%%%%
%% optional
%\supplementary{The following supporting information can be downloaded at:  \linksupplementary{s1}, Figure S1: title; Table S1: title; Video S1: title.}

%%%%%%%%%%%%%%%%%%%%%%%%%%%%%%%%%%%%%%%%%%
\authorcontributions{
Conceptualization, M.M.; methodology, M.M.; formal analysis, M.M.; investigation, M.M.; writing---original draft preparation, M.M.; writing---review and editing, M.M. The author has read and agreed to the published version of the manuscript.}

\funding{This research received no external funding.}

\acknowledgments{The author thanks colleagues at the Faculty of Informatics, Juraj Dobrila University of Pula, for discussions that informed the framework.}

\dataavailability{No new data were created or analyzed in this study. Data sharing is not applicable to this article.}

\conflictsofinterest{The author declares no conflicts of interest.} 

%%%%%%%%%%%%%%%%%%%%%%%%%%%%%%%%%%%%%%%%%%
%% Optional

%% Only for journal Encyclopedia
%\entrylink{The Link to this entry published on the encyclopedia platform.}

\abbreviations{Abbreviations}{
The following abbreviations are used in this manuscript:

\noindent
\begin{tabular}{@{}ll}
AI    & Artificial Intelligence \\
ABC   & Arena Blended Connected (Learning Design) \\
KABL-ID & Cognitive Alignment Baseline for Instructional Design \\
PRISMA & Preferred Reporting Items for Systematic Reviews and Meta-Analyses \\
LD    & Learning Design \\
\end{tabular}
}

%%%%%%%%%%%%%%%%%%%%%%%%%%%%%%%%%%%%%%%%%%
\appendixtitles{no}
\appendixstart
\appendix

\section[\appendixname~\thesection]{PRISMA Checklist Summary}
\label{appendix:prisma}

\begin{table}[H]
\caption{PRISMA 2020 checklist items addressed in this paper (abbreviated).}
\centering
\small
\begin{tabularx}{\linewidth}{lX}
\toprule
Item & Location / Status \\
\midrule
Title identifies as systematic review & Partially (combined with conceptual framework) \\
Rationale & Section 1 \\
Objectives & Section 1 \\
Eligibility criteria & Section 3.2.2 \\
Information sources & Section 3.2.3 \\
Search strategy & Section 3.2.3 \\
Selection process & Section 3.2.4 \\
Data collection process & Section 3.2.5 \\
Data items & Section 3.2.5 \\
Study risk of bias & Not formally assessed; acknowledged as limitation \\
Synthesis methods & Narrative synthesis; Section 4.1 \\
Study selection results & Section 4.1 \\
Study characteristics & Section 4.1 \\
Results of syntheses & Section 4.1 \\
Limitations of evidence & Section 6 \\
Interpretation & Section 5 \\
Funding & Funding statement \\
Registration & Not registered; acknowledged \\
\bottomrule
\end{tabularx}
\end{table}

\section[\appendixname~\thesection]{Curriculum Coding Record}
\label{appendix:coding}

The coding record for the four outcomes in Table~\ref{tab:curriculum} follows the KABL-ID protocol: each outcome was decomposed into its action verb, the Bloom level was assigned according to the revised taxonomy verb lists \cite{anderson_krathwohl_2001}, and the ABC type was assigned based on the activity structure described in the outcome.

For a full coding of all outcomes in the Computational Thinking and Programming domain across all grade levels, the same protocol applies. The author's working coding sheet is available on request.

%%%%%%%%%%%%%%%%%%%%%%%%%%%%%%%%%%%%%%%%%%
%\isPreprints{}{% This command is only used for ``preprints''.
\begin{adjustwidth}{-\extralength}{0cm}
%} % If the paper is ``preprints'', please uncomment this parenthesis.
%\printendnotes[custom] % Un-comment to print a list of endnotes

\reftitle{References}

\bibliography{references}

\PublishersNote{}
%\isPreprints{}{% This command is only used for ``preprints''.
\end{adjustwidth}
%} % If the paper is ``preprints'', please uncomment this parenthesis.
\end{document}